% =========================================================
% config/environments.tex
% Environment khusus buku ajar berbasis OBE.
% =========================================================

% ---------------------------------------------------------
% Penambal kompatibilitas: enumitem + tcolorbox breakable.
%
% Package enumitem menetapkan \enit@after secara LOKAL pada saat
% \begin{enumerate}. Bila sebuah daftar berada di dalam tcolorbox
% yang terpotong antar halaman, tcolorbox menyusun kedua potongan
% di dalam grup yang berbeda, sehingga \end{enumerate} pada
% potongan kedua menemukan \enit@after dalam keadaan tak
% terdefinisi dan perakitan gagal.
%
% Menetapkan nilai baku secara global menghilangkan kegagalan itu
% tanpa mengubah perilaku daftar, karena \@empty memang nilai baku
% \enit@after.
% ---------------------------------------------------------
\makeatletter
\AtBeginDocument{%
    \@ifundefined{enit@after}{\let\enit@after\@empty}{}%
}
\makeatother

% ---------------------------------------------------------
% Teorema-seperti (amsthm)
% ---------------------------------------------------------
\theoremstyle{definition}
\newtheorem{definisi}{Definisi}[chapter]
\newtheorem{contoh}{Contoh}[chapter]
\newtheorem{prosedur}{Prosedur}[chapter]

\theoremstyle{plain}
\newtheorem{teorema}{Teorema}[chapter]
\newtheorem{lemma}{Lemma}[chapter]

\theoremstyle{remark}
\newtheorem{catatan}{Catatan}[chapter]

% Catatan: palet warna buku didefinisikan di config/styles.tex,
% yang dimuat lebih dahulu melalui preamble.tex.

% ---------------------------------------------------------
% Kotak OBE
% Semua kotak dapat terpotong antar halaman (breakable).
% ---------------------------------------------------------
\tcbset{
    obe box/.style={
        breakable,
        enhanced,
        sharp corners,
        boxrule=0pt,
        leftrule=3pt,
        left=6pt,
        right=6pt,
        top=6pt,
        bottom=6pt,
        fonttitle=\bfseries,
        coltitle=white,
        attach boxed title to top left={
            xshift=6pt,
            yshift=-\tcboxedtitleheight/2
        },
        boxed title style={
            sharp corners,
            boxrule=0pt,
            left=5pt,
            right=5pt,
            top=2pt,
            bottom=2pt
        }
    }
}

% --- Sub-CPMK bab ---
\newtcolorbox{subcpmkbox}[1][]{
    obe box,
    colback=obeBlueLight,
    colframe=obeBlue,
    colbacktitle=obeBlue,
    title={Outcome Pembelajaran Khusus Bab (Sub-CPMK)},
    #1
}

% --- Tujuan pembelajaran ---
\newtcolorbox{tujuanbox}[1][]{
    obe box,
    colback=obeBlueLight,
    colframe=obeBlue,
    colbacktitle=obeBlue,
    title={Tujuan Pembelajaran},
    #1
}

% --- Materi pokok ringkas ---
\newtcolorbox{materibox}[1][]{
    obe box,
    colback=obeGrayLight,
    colframe=obeGray,
    colbacktitle=obeGray,
    title={Materi Pokok},
    #1
}

% --- Aktivitas pembelajaran ---
\newtcolorbox{aktivitasbox}[1][]{
    obe box,
    colback=obeTealLight,
    colframe=obeTeal,
    colbacktitle=obeTeal,
    title={Aktivitas Pembelajaran yang Direkomendasikan},
    #1
}

% --- Praktikum ---
\newtcolorbox{praktikumbox}[1][]{
    obe box,
    colback=obeTealLight,
    colframe=obeTeal,
    colbacktitle=obeTeal,
    title={Praktikum},
    #1
}

% --- Asesmen ---
\newtcolorbox{asesmenbox}[1][]{
    obe box,
    colback=obeAmberLight,
    colframe=obeAmber,
    colbacktitle=obeAmber,
    title={Asesmen (Evaluasi Kinerja)},
    #1
}

% --- Checklist kompetensi ---
\newtcolorbox{checklistbox}[1][]{
    obe box,
    colback=obeGreenLight,
    colframe=obeGreen,
    colbacktitle=obeGreen,
    title={Checklist Pencapaian Kompetensi},
    #1
}

% --- Refleksi ---
\newtcolorbox{refleksibox}[1][]{
    obe box,
    colback=obeGreenLight,
    colframe=obeGreen,
    colbacktitle=obeGreen,
    title={Refleksi},
    #1
}

% --- Studi kasus ---
\newtcolorbox{studikasusbox}[1][]{
    obe box,
    colback=obeMaroonLight,
    colframe=obeMaroon,
    colbacktitle=obeMaroon,
    title={Studi Kasus},
    #1
}

% --- Peringatan / kesalahan umum ---
\newtcolorbox{peringatanbox}[1][]{
    obe box,
    colback=obeMaroonLight,
    colframe=obeMaroon,
    colbacktitle=obeMaroon,
    title={Kesalahan yang Sering Terjadi},
    #1
}

% --- Catatan tambahan ---
\newtcolorbox{infobox}[1][]{
    obe box,
    colback=obeGrayLight,
    colframe=obeGray,
    colbacktitle=obeGray,
    title={Catatan},
    #1
}

% --- Ringkasan instruksi ---
\newtcolorbox{instruksibox}[1][]{
    obe box,
    colback=obeBlueLight,
    colframe=obeBlue,
    colbacktitle=obeBlue,
    title={Ringkasan Instruksi},
    #1
}
